\documentclass[twocolumn,superscriptaddress,aps,prl,floatfix]{revtex4-2}

% ============================================================
% DiffCMB — Draft manuscript
% "A full-sky differentiable joint sampler over the CMB lensing potential"
%
% DRAFTING RULES (ROADMAP.md, 2026-08-06):
%   - NEVER lead with the differentiable-SHT machinery. Flinch (arXiv:2510.26691)
%     made curved-sky differentiable inference table stakes. Every abstract, intro
%     sentence and talk title leads with the JOINT (C_ell, C_L^phiphi) POSTERIOR.
%   - Differentiable infrastructure goes in Methods as "enabling machinery,
%     consistent with Flinch", never in the novelty sentence.
%   - Every "first" carries scope qualifiers and nearest-prior-work citations.
%   - Re-check the named-author arXiv scan (Millea, Seljak, Bayer, Loureiro)
%     before every submission milestone.
%
% SECTION STATUS:
%   [DONE]        = content fixed, does not depend on pending results
%   [PLACEHOLDER] = leave stub; will fill once chains are done
%   [TODO]        = not started; needed before submission
% ============================================================

\usepackage{amsmath,amssymb,bm}
\usepackage{graphicx}
\usepackage{hyperref}
\usepackage{xcolor}
\usepackage{natbib}

% Convenience macros
\newcommand{\alm}{a_{\ell m}}
\newcommand{\cl}{C_\ell}
\newcommand{\clpp}{C_L^{\phi\phi}}
\newcommand{\phiphi}{\phi\phi}
\newcommand{\lmax}{\ell_{\max}}
\newcommand{\todo}[1]{\textcolor{red}{\textbf{[TODO: #1]}}}
\newcommand{\placeholder}[1]{\textcolor{blue}{\textbf{[PLACEHOLDER: #1]}}}

% ============================================================
\begin{document}

\title{A full-sky differentiable joint Gibbs sampler over
       the CMB temperature field, power spectrum, and lensing potential}

\author{\placeholder{Oscar Hickman}}
\affiliation{\placeholder{Institute for Computational Cosmology, Durham University}}

% Additional authors to be confirmed
\date{\today}

% ============================================================
\begin{abstract}
% RULE: lead with the joint posterior object, not the machinery.
% Draft language from ROADMAP.md 2026-08-06 re-pitch:
We present the first curved-sky, full-sky posterior over the joint
$(\alm, \cl, \phi, \clpp)$ — a four-block Gibbs sampler operating directly on
HEALPix maps whose blocks alternate exact inverse-Gamma draws for the two power
spectra and Hamiltonian Monte Carlo / microcanonical Langevin steps for the two
field amplitudes.
The sampler's defining output is the joint $(\cl^{TT}, \clpp)$ posterior with
correlations propagated rather than assumed, an object no existing method
produces: quadratic estimators and iterative MAP reconstruction (e.g.\
\citealt{Carron:2017, MUSE}) return point estimates or marginal constraints;
Commander-class samplers \citep{Commander} and their differentiable descendants
\citep{Almanac,Flinch} infer maps and power spectra on the full sky but carry no
lensing-potential block.
We validate the sampler via simulation-based calibration (SBC) on lensed
temperature maps at $\lmax = \placeholder{64}$, demonstrating frequentist
coverage of the joint posterior.
As an exact reference standard, the sampler's posterior is the natural
calibration target for fast, marginal, or learned CMB lensing methods.
\todo{Add key numbers once chains are done: coverage pass, joint-correlation amplitude.}
\end{abstract}

\maketitle

% ============================================================
\section{Introduction}
\label{sec:intro}
% STATUS: [DONE — draft from ROADMAP.md re-pitch; sharpen at writing time]

% --- Opening: what the competing landscape is ---
Curved-sky differentiable field-level CMB inference is now established.
\citet{Flinch} demonstrate gradient propagation from masked full-sky temperature
maps to cosmological parameters, recovering maps \textit{and} angular power
spectra;
\citet{Almanac} sample all-sky CMB maps jointly with their power spectra via
Hamiltonian Monte Carlo.
Both analyses are \textit{lensing-blind}: they infer $(\alm, \cl)$ and, in
Flinch's case, cosmological parameters, but carry no lensing-potential $\phi$ and
no $\clpp$ block.

Conversely, the methods that do treat lensing on the curved sky — quadratic
estimators, iterative maximum \textit{a posteriori} (MAP) reconstruction
\citep{Carron:2017}, and MUSE \citep{MUSE} — return point estimates or marginal
constraints, not a joint posterior.
The flat-sky code CMBLensing.jl \citep{CMBLensing} samples jointly over
$(\alm, \phi)$ on the flat sky but does not cover the full sky.

This paper closes the remaining cell: a full-sky, HEALPix, differentiable joint
Gibbs sampler over $(\alm, \cl, \phi, \clpp)$, validated by
simulation-based calibration.
Its primary output is the joint $(\cl^{TT}, \clpp)$ posterior and its
correlations — an object no existing method produces — which functions as a
\textit{reference standard}: the exact posterior that fast, marginal, or learned
methods \citep{Doeser:2024} get validated against.

% --- What we offer vs each competing paradigm ---
\subsection*{Position vs existing methods}

\textbf{vs Commander / Almanac / Flinch} (exact field samplers, lensing-blind):
our sampler adds the $(\phi, \clpp)$ blocks.
The differentiable spherical-harmonic infrastructure we use is consistent with
\citet{Flinch} and treated here as enabling machinery, not as a contribution.

\textbf{vs QE / iterative MAP / MUSE} (lensing-sensitive, marginal):
our sampler provides a joint posterior rather than a point estimate or a
marginalised constraint; it is not competitive in speed but provides
the calibration target these methods lack.

\textbf{vs diffusion / score-based generative methods} \citep{Drouin:2023}:
diffusion samplers produce approximate uncorrelated samples quickly but discard
the two distinguishing properties we build: a differentiable forward model and
a statistically exact sampler.
An exact sampler is what learned posteriors get validated against
\citep{Doeser:2024}.
The comparison is not speed but correctness.

\textbf{Honest scope statement.}
We demonstrate coverage at $\lmax = \placeholder{64}$. We quote the scale at
which the calibration test actually passes rather than the largest scale at
which a chain has been run: a converged result at a smaller scale is worth more
than a partially-converged one at a larger scale, and the two should not be
conflated.
Scaling to $\lmax \sim 1000$ and polarisation (the science-driver configuration)
is left to follow-up work.


% ============================================================
\section{Method}
\label{sec:method}
% STATUS: [DONE for block structure; sub-sections to be fleshed out]

We describe the four-block Gibbs sampler.
Let $\mathbf{d}$ denote the observed (lensed, beam-convolved, noise-added)
temperature map.
The joint target is
\begin{equation}
    p(\alm, \cl, \phi, \clpp \mid \mathbf{d})
    \propto p(\mathbf{d} \mid \alm, \phi)\,
            p(\alm \mid \cl)\,p(\cl)\,p(\phi \mid \clpp)\,p(\clpp),
\end{equation}
where $p(\mathbf{d} \mid \alm, \phi)$ is the lensed CMB likelihood evaluated via
a differentiable full-sky lensing operator.

% --- Block 1 ---
\subsection{Block 1: $\cl \mid \alm$ (exact inverse-Gamma draw)}
\label{sec:block1}

Given the current unlensed alm amplitudes, the conditional $p(\cl \mid \alm)$
is exactly inverse-Gamma (conjugate to the Gaussian alm prior) and is sampled
exactly by a one-line draw.
This is the standard Gibbs move used in Commander \citep{Commander} and
Almanac \citep{Almanac}.

% --- Block 2 ---
\subsection{Block 2: $\alm \mid \cl, \phi$ (HMC)}
\label{sec:block2}

We sample the unlensed alm given $\cl$ and $\phi$ via Hamiltonian Monte Carlo,
using the matrix-free ducc0 spherical-harmonic transform \citep{Reinecke:2023}
wrapped in a \texttt{tf.custom\_gradient} to propagate autodiff gradients through
the lensing operator.
This is enabling machinery consistent with \citet{Flinch}; the differentiable
SHT is not presented as a contribution.
The mass matrix is the posterior precision $1/\cl$; whitened coordinates are
used throughout.

% --- Block 3 ---
\subsection{Block 3: $\phi \mid \alm, \cl$ (HMC)}
\label{sec:block3}

We sample the lensing potential $\phi$ given the current $\alm$ and $\cl$ with
Hamiltonian Monte Carlo, using the differentiable lensed log-posterior of
Sec.~\ref{sec:lensing} and its autodiff gradient, a prior-preconditioned
diagonal mass matrix, and a long trajectory
(\placeholder{$n_{\rm leapfrog} = 240$}) --- the trajectory length, rather than
the step size, is what controls $\phi$ mixing in this target, at an acceptance
rate of \placeholder{0.7\text{--}0.8}.

We record the alternatives tried, since the literature would suggest otherwise.
Microcanonical Langevin Monte Carlo \citep{Robnik:2022} is reported by
\citet{Bayer:2023,Flinch} to give one-to-three-order-of-magnitude efficiency
gains over HMC at this dimensionality on curved-sky CMB fields, and we
implemented it against the same target. It wins on effective samples per
wall-clock second at every multipole bin we probed, but fails our stationarity
gate in both an untuned and a step-size-tuned configuration --- in one tuned run
the lag-1 autocorrelation was \emph{worse} than untuned. The No-U-Turn Sampler
and a Fisher-based mass matrix were likewise tested and rejected. We therefore
use plain HMC, and note that an efficiency metric which ignores whether the
chain is stationary will select the wrong sampler here.

% --- Block 4 ---
\subsection{Block 4: $\clpp \mid \phi$ (exact inverse-Gamma draw)}
\label{sec:block4}

Given the current $\phi$ amplitudes, the conditional $p(\clpp \mid \phi)$ is
exactly inverse-Gamma by the same conjugacy argument as Block 1,
\begin{equation}
  \clpp \mid \phi \;\sim\; \mathrm{InvGamma}\!\left(k_L/2 - 1,\; S_L/2\right),
  \qquad S_L = \sum_{m=-L}^{L} |\phi_{Lm}|^2 ,
  \label{eq:block4}
\end{equation}
and this block resamples the $\phi$ prior spectrum every sweep and rebuilds the
Block 3 mass matrix to match, closing the joint posterior.

\subsubsection*{The prior on \texorpdfstring{$\clpp$}{ClPhiPhi}, and why it is
not a free choice}
\label{sec:clpp-prior}

Equation~\eqref{eq:block4} deserves care in two respects: the number of degrees
of freedom it encodes, and the prior it implies.

Our packed parameterisation stores an imaginary part only for $m \ge 2$, so
each multipole carries $k_L = 2L$ real degrees of freedom rather than $2L+1$.
The Gaussian prior normalisation is then $C^{-k_L/2}$, and against a prior
$\mathrm{InvGamma}(a_0, b_0)$ the exact conditional is
$\mathrm{InvGamma}(k_L/2 + a_0,\; b_0 + S_L/2)$ --- i.e.\ shape $L-1$ under the
flat prior, not $L-\tfrac12$. We note this explicitly because the error is
invisible to the obvious check: a probability-integral-transform test of the
draws against the sampler's own stated conditional passes for \emph{any} shape
parameter, since code and reference share the assumption. Only an independent
generative calibration test --- which requires a proper prior, below --- can
detect it.

Equation~\eqref{eq:block4} is otherwise the conditional implied by a
\emph{flat} prior on $\clpp$, which is improper. This is worth stating explicitly rather than
leaving implicit, because it propagates into the $\phi$ block in a way that is
easy to miss. Integrating $\clpp$ out of the joint gives a marginal on the
field, $p(\phi) \propto S_L^{-(k_L/2 - 1)} = S_L^{-(L-1)}$; against the radial
measure for the $k_L = 2L$ real components of $\phi_{Lm}$,
$r^{\,k_L - 1}\,\mathrm{d}r$ with $S_L = r^2$, this is
\begin{equation}
  p(r) \;\propto\; r^{\,1} ,
\end{equation}
i.e.\ flat in $S_L$ --- improper, and \emph{increasing} with amplitude. The
$\phi$ amplitude is therefore constrained by the lensing likelihood alone, with
no prior restoring force. Equivalently, substituting the conditional mean of
Eq.~\eqref{eq:block4} into the Gaussian $\phi$ prior term gives
$\tfrac12 \sum_L S_L/\clpp = \sum_L (L - 2)$, independent of $\phi$: the
prior is exactly scale-free along the amplitude ray.

In practice this is benign whenever the lensing signal-to-noise is sufficient
to pin the amplitude, which it is in the configurations reported here. It is
not benign as a matter of principle, and it has two concrete consequences that
we report rather than gloss: it removes the only safeguard against an
amplitude runaway if the temperature block is poorly initialised, and it means
a rank test of $\phi$ against a truth drawn from a \emph{proper} Gaussian prior
is not a calibration test of this target (Sec.~\ref{sec:sbc}).

We therefore also provide the conjugate proper alternative, a weakly
informative $\mathrm{InvGamma}(\nu/2,\, \nu C_L^{\phi\phi,\,\rm fid}/2)$ prior centred on
a fiducial spectrum, under which
\begin{equation}
  \clpp \mid \phi \;\sim\; \mathrm{InvGamma}\!\left(k_L/2 + \tfrac{\nu}{2},\;
  \big(S_L + \nu\,C_L^{\phi\phi,\,\rm fid}\big)/2\right).
  \label{eq:block4-proper}
\end{equation}
Here $\nu$ reads as an effective number of prior pseudo-modes, on the same
footing as the $k_L$ real modes the data supplies at multipole $L$, so the
prior regularises the lowest multipoles --- where the improper prior does its
damage --- and is negligible at high $L$. The marginal tail becomes
$p(r) \propto r^{\,-1-\nu}$, so the target is proper exactly for $\nu > 0$; we
use $\nu = \placeholder{6}$, and reject $\nu \le 0$ at the API level rather
than allowing a choice that appears to fix the problem while leaving it in
place.

% --- Forward model ---
\subsection{Differentiable lensing operator}
\label{sec:lensing}

The lensing forward model maps $(\alm, \phi) \to T_{\rm lensed}$ via bilinear
interpolation on the sphere, validated against the \texttt{healpy} deflection
reference and the dense Jacobian of \citet{Reinecke:2023} to machine precision.
Gradients with respect to both $\alm$ and $\phi$ are available via
\texttt{tf.custom\_gradient} / \texttt{tf.py\_function} wrapping of the ducc0
SHT; this is the same approach used in \citet{Flinch} and is treated here
as a standard building block.

% --- What the bias is NOT ---
\subsection{What the lensing-potential deficit is not}
\label{sec:deficit-not}

% ROADMAP.md: "Ruled out by construction — say so in three sentences in the
% paper, don't spend compute."
Several possible sources of systematic bias in the recovered $\phi$ posterior
are ruled out by the construction of this sampler and need not be tested
empirically.
(i)~\textit{N0/N1 noise bias}: a Gibbs sampler has no noise-bias subtraction
step; any bias in the recovered $\phi$ amplitude cannot originate from the
multiplicative noise floor of a quadratic estimator.
(ii)~\textit{Non-Gaussian deflections}: our validation simulations draw $\phi$
from its Gaussian prior exactly; non-Gaussianity of the lensing field is not a
factor.
(iii)~\textit{Foreground contamination}: the validation uses foreground-free
CMB-only simulations.
(iv)~\textit{Lensing-operator accuracy}: the matrix-free SHT forward model is
validated against the dense reference and independent \texttt{healpy} outputs to
numerical precision, following the standard set by \citet{Reinecke:2023}.
(v)~\textit{Chain initialisation}: an earlier version of these ensembles began
the $\alm$ block from a cold prior draw rather than the data-driven warm start,
and the $\phi$ amplitude then ran away by orders of magnitude, since with a
resampled $\clpp$ nothing in the target restrains it (Sec.~\ref{sec:clpp-prior})
and $\phi$ is the only freedom left to absorb the $\alm$ residual. The effect
was diagnostic rather than subtle --- across realizations the $\phi$ inflation
correlated with the $\alm$ misfit at $\rho = \placeholder{-0.78}$
($p = \placeholder{0.003}$) --- and all results reported here use the warm
start, with $\phi$ power within a factor $\placeholder{0.8\text{--}1.4}$ of the
truth in every chain.


% ============================================================
\section{Validation}
\label{sec:results}
% STATUS: [PLACEHOLDER — content depends on the equilibrated chains]

\subsection{Simulation-based calibration (SBC)}
\label{sec:sbc}

We run 12 independent chains at $\lmax = \placeholder{64}$, each with its own
$(\alm^{\rm true}, \phi^{\rm true}, \mathrm{noise})$ triple drawn from the
fiducial spectra and warm-started away from the truth, and rank the truth's
binned power within each posterior.

\subsubsection*{The rank statistic requires its own null}

The spectrum-level statistic must be calibrated before any departure from
uniformity is read as bias, and this is not a technicality. Ranking the truth's
\emph{realized} power $S_L/k_L$ against posterior draws of $\clpp$ compares
it against the \emph{mode} of Eq.~\eqref{eq:block4}, since
$\beta/(\alpha+1) = S_L/k_L$. The inverse-Gamma is right-skewed, so
$P(\text{draw} < \text{mode}) < 1/2$ \emph{for an exact sampler}, and averaging
over $\ell$ within a bin shrinks the spread while preserving the offset, driving
the mean rank toward zero. A naive uniformity test therefore rejects a correct
sampler with arbitrarily high confidence at large bin width.

We accordingly compare each observed rank against a null simulated from the
conditional itself. For $\cl^{TT}$ the observed bin means
$\placeholder{0.083, 0.083, 0.094, 0.333}$ agree with the simulated null
$\placeholder{0.086, 0.092, 0.114, 0.344}$ in every bin, each observation
lying inside the null's $95\%$ band. For $\clpp$ the null
must additionally retain the chain's own sweep-to-sweep scatter in $\phi$,
which contributes a common-mode variance across $\ell$ that does not average
down within a bin; a null that freezes $\phi$ at the truth is far too narrow
and makes a correct sampler look biased. With that included, the observed
$\placeholder{0.260, 0.302, 0.427, 0.406}$ again lies within the null band
$\placeholder{0.279, 0.339, 0.457, 0.376}$ throughout. We report these rows as interval coverage, not calibration, since
the implied $\cl$ prior is flat and improper and no strict rank exists.

\subsubsection*{Block-level exactness}

Independently of the rank test, we verify Block 4 directly on the production
chains: using each sweep's $(\phi, \clpp)$ pair, the probability integral
transform $u = F_{\mathrm{InvGamma}(k_L/2 + a_0,\,(S_L(\phi)+b_{0,L})/2)}(\clpp)$
is uniform (\placeholder{$\mathrm{KS}\ p = 0.42$}, \placeholder{$\bar u = 0.4999$},
\placeholder{$N = 17856$}). A probability-integral transform is only evidence
if a deliberately misaligned pairing is rejected, and the lag must exceed the
scale on which the conditioning variable actually moves: at lag~1 the chain's
$S_L(\phi)$ autocorrelation is still \placeholder{0.92} and the control is
\emph{not} rejected (\placeholder{$p = 0.18$}), whereas at lag
\placeholder{10} and \placeholder{50} (autocorrelation \placeholder{0.52} and
\placeholder{0.09}) it is rejected with \placeholder{$p < 10^{-16}$}. The
aligned pass is therefore a genuine one: Block~4 draws from its stated
conditional given $\phi$. We stress that this validates the \emph{draw}, not
the derivation --- code and reference share the same expression, so this test
passes for any shape parameter, and it is precisely how the incorrect $k_L$ of
Sec.~\ref{sec:clpp-prior} survived an earlier check. The derivation is tested
by the independent generative rank below.

\subsubsection*{Field-level calibration}

The field ranks are the statistic this ensemble design supports directly, and
they are reported for the configuration in which they are interpretable: with
$\clpp$ held at the fiducial spectrum, so that the $\phi$ prior is proper and
\emph{identical to the process that generated the truth}. Pooling 12
independent chains at $\lmax = \placeholder{64}$ over four multipole bins, the
mean normalised ranks are
\begin{equation*}
  \bar u_{\phi} = \placeholder{0.4688}\ (p = \placeholder{0.124}), \qquad
  \bar u_{\alm} = \placeholder{0.5312}\ (p = \placeholder{0.235}),
\end{equation*}
both consistent with uniformity and with no individual bin flagged in either
field, against $\bar u_\phi = \placeholder{0.367}$
($p = \placeholder{0.004}$) for the same configuration with $\clpp$ resampled
under the flat improper prior. These are measured against the corrected
inverse-Gamma shape of Sec.~\ref{sec:clpp-prior}; the configuration is
otherwise a single-variable change from the resampled-$\clpp$ run, so the
difference is attributable to the $\clpp$ prior alone.
\todo{Add rank histograms.}

That the $\phi$ rank is uniform only once $\clpp$ is held fixed is itself
informative, and is the reason we treat the prior of Sec.~\ref{sec:clpp-prior}
as a modelling decision rather than a detail. Under the flat prior the target's
$\phi$ marginal is flat in $S_L$ while $\phi^{\rm true}$ is drawn from a proper
Gaussian, so the two distributions differ by construction and the rank has no
reason to be uniform: the departure is a statement about the prior, not
evidence of a sampler defect. Adopting the proper conjugate prior of
Eq.~\eqref{eq:block4-proper} and drawing the truth from that same joint prior
restores a strict calibration test while retaining the $\clpp$ block that
produces the joint posterior of Sec.~\ref{sec:joint-posterior}.

We report the status of that configuration rather than assert it. Run with
$\nu = \placeholder{6}$ at $\lmax = \placeholder{64}$, the strict rank of
the true $\clpp$ --- drawn, uniquely in this configuration, from the same
prior the sampler targets, so that no null simulation is required --- is
$\bar u = \placeholder{0.3802}$ ($p = \placeholder{0.0013}$), against
$\placeholder{0.25\text{--}0.28}$ ($p < 10^{-4}$) before the shape correction
of Sec.~\ref{sec:clpp-prior}. The $\alm$ rank is clean
($\placeholder{0.5052}$, $p = \placeholder{0.41}$) and Block~4 is exact given
$\phi$, so the residual is localised to the $\phi$ block or its mixing under
the funnel that resampling $\clpp$ induces: the same chains mix markedly worse
than the fixed-$\clpp$ ensemble ($\tau_{\rm int}$ median
$\placeholder{24\text{--}56}$ versus $\placeholder{4.7\text{--}27}$ per bin),
and the deficit is not removed by discarding the first half or three quarters
of each chain. We therefore quote the calibration claim from the fixed-$\clpp$
configuration above, and present the $\clpp$ posterior of
Sec.~\ref{sec:joint-posterior} as the posterior this sampler produces rather
than as a calibration-validated one.
\todo{Resolve with a longer-trajectory rerun; update if it moves to uniform.}

\subsection{Joint $(\cl^{TT}, \clpp)$ posterior}
\label{sec:joint-posterior}

\placeholder{
Figure: 2D marginal $(\cl^{TT}, \clpp)$ posterior from the pooled coverage
chains, showing the cross-correlation between the two power spectra.
This is the headline result: no existing method produces this joint object.
Result: \todo{correlation amplitude and sign.}
}

This figure is sourced from the proper-prior configuration of
Eq.~\eqref{eq:block4-proper}, which is the only one in which $\clpp$ is a
proper Bayesian object (Sec.~\ref{sec:clpp-prior}). It therefore carries the
caveat established in Sec.~\ref{sec:sbc}: that configuration's strict $\clpp$
rank is $\placeholder{0.3802}$ ($p = \placeholder{0.0013}$), improved but not
yet uniform, so the joint posterior is presented as the object this sampler
produces and not as a calibration-validated one.

\subsection{C$_\ell^{TT}$ bias vs lensing-blind analysis}
\label{sec:cl-bias}

\placeholder{
Figure: recovered $\cl^{TT}$ from the lensing-aware joint sampler vs the
Commander-style lensing-blind Gibbs (same simulation, same noise, no $\phi$
block). The lensing-blind chain fits the lensed-sky power spectrum as if it were
unlensed; the bias shows directly how much lensing power leaks into the
inferred $\cl^{TT}$.
Data already available: \texttt{results/analysis/lensing\_blind\_baseline\_lmax128.npz}
(run submitted 2026-08-10, job 11717687).
}

\subsection{Per-mode uncertainty propagation}
\label{sec:uncertainty}

\placeholder{
Figure: posterior width on $\phi_{\ell m}$ from the joint sampler vs the
posterior width one would assign under the marginal (QE) approximation.
Quantifies what joint sampling buys over a marginalised method.
}


\subsection{Comparison with the flat-sky joint sampler CMBLensing.jl}
\label{sec:cmblensing}

CMBLensing.jl \citep{CMBLensing,Millea:2020} is the closest existing method:
the only other code that samples a joint $(\text{field}, \phi)$ posterior rather
than marginalising or point-estimating. It is flat-sky, which is the gap this
work closes, so the comparison is a positioning statement rather than a
head-to-head benchmark, and we make it from that code's published
configurations rather than by re-running it --- their paper reports the
configuration table, autocorrelation lengths and wall-times needed for every
axis compared here.

\emph{What matches.} Our patch-scale validation used a polar-cap mask at
$f_{\rm sky} = \placeholder{0.016}$, or $\approx \placeholder{660}\,\mathrm{deg}^2$,
which is within a few percent of their largest published configuration
(\textsc{big}, $\placeholder{650}\,\mathrm{deg}^2$) --- a coincidence of
independent design choices rather than a matched setup, but it means the two
codes are being compared at the same sky area. That configuration is described
by its authors as near the practical ceiling for a flat-sky treatment, which is
the point: the curved-sky formulation is not an alternative implementation of
the same calculation, it is what removes that ceiling.

\emph{What does not match, stated rather than smoothed over.} Their published
examples are polarisation-only ($QU$) while ours is temperature-only; their
Fourier mask reaches $\ell < \placeholder{3500}$ while our validated production
scale is $\lmax = \placeholder{300}$, an order of magnitude in $\ell$; and the
two codes' fiducial cosmologies differ in $\tau$, $\omega_b$ and $\omega_c$.
Each of these affects lensing sensitivity and the degeneracy structure, so the
comparison is one of \emph{methodology and mixing behaviour}, not of information
content. We state it that way rather than inflating our configuration for the
comparison.

\emph{Architecture.} Both codes are Gibbs samplers with one non-Gaussian block
on HMC, and the resemblance is close enough that one difference needs
disambiguating. CMBLensing.jl draws its field block by conjugate gradient; we
tested a conjugate-gradient shortcut for our $\alm$ block and rejected it as
biased once a $\phi$ block was active. These are not the same thing. Their
field-block conditional is genuinely Gaussian given fixed $\phi$, with $\phi$
drawn in a separate HMC pass --- the same architecture we use --- whereas the
route we rejected replaced the joint draw itself. A reader who sees ``they use
CG, they rejected CG'' should read it as a difference in what was conditioned
on, not a contradiction.

\emph{Cost.} Their three configurations report
\placeholder{19\text{--}50}\,h on a single GPU with per-parameter
autocorrelation lengths of \placeholder{5\text{--}33}; ours runs at
\placeholder{29\text{--}54}\,s per Gibbs sweep at $\lmax = \placeholder{300}$
with $\phi$-block effective sample fractions of
\placeholder{4.7\text{--}11.7\%}. We deliberately do not reduce these to a
single cost-per-effective-sample ratio: the two are measured on different
observables, at different $\ell$ ranges, on different hardware, and a single
number would imply a precision the comparison does not have.

\todo{If a referee asks for same-realization posterior overlays, that is the
trigger to install and run CMBLensing.jl --- published numbers cannot provide
them. Not before.}

% ============================================================
\section{Discussion and conclusion}
\label{sec:discussion}
% STATUS: [PARTIAL — framing is fixed; numbers are placeholders]

We have presented the first full-sky, curved-sky joint Gibbs sampler over
$(\alm, \cl, \phi, \clpp)$.
The sampler's defining output — the joint $(\cl^{TT}, \clpp)$ posterior with
propagated correlations — is not produced by any existing method:
Commander-class samplers \citep{Commander,Almanac,Flinch} and QE-based methods
\citep{Carron:2017} both produce either the map-and-spectra posterior
(lensing-blind) or a lensing constraint (marginalised), never both together.

\textbf{Honest scale statement.}
The demonstration is at $\lmax = \placeholder{64}$, where chains
equilibrate in $\sim \placeholder{XX}$h on a single GPU node.
Scaling to $\lmax \sim 1000$ (the Planck / SO lensing science scale) is the
principal avenue for follow-up work and is not addressed here.

\textbf{As a reference standard.}
A sample-exact posterior is the object that amortised or learned CMB lensing
posteriors \citep{Doeser:2024} ultimately get validated against.
We offer the joint $(\cl^{TT}, \clpp)$ posterior at demonstration scale as
that reference.

\textbf{Outlook.}
Immediate extensions include:
real Planck data (\placeholder{in prep});
$\Lambda$CDM parameter inference from the posterior $\cl^{TT}$ chains;
polarisation (TQU) and LiteBIRD delensing forecasts via the spin-2 SHT
extension of the current infrastructure.

% ============================================================
\section*{Acknowledgements}
\todo{Acknowledgements to be added.}

% ============================================================
% REFERENCES
% Use inline \bibitem entries until a proper .bib file is assembled.
\begin{thebibliography}{99}

\bibitem{Flinch}
Crespi, S., Bonici, M., Loureiro, A., et al.\ 2025,
\textit{Flinch: A Differentiable Framework for Field-Level Inference of Cosmological Parameters from Curved Sky Data},
arXiv:2510.26691

\bibitem{Almanac}
Sellentin, E., Loureiro, A., Whiteway, L., et al.\ 2023,
\textit{Almanac: Scalable all-sky HMC CMB analysis},
arXiv:2305.16134

\bibitem{Commander}
Eriksen, H.\ K., et al.\ 2004,
\textit{Power Spectrum Estimation from the WMAP Data Using the Master HEALPix Estimator},
ApJS, 155, 227

\bibitem{Carron:2017}
Carron, J.\ \& Lewis, A.\ 2017,
\textit{Maximum a posteriori CMB lensing reconstruction},
PRD, 96, 063510,
arXiv:1704.08230

\bibitem{MUSE}
Millea, M.\ \& Seljak, U.\ 2022,
\textit{Marginal Unbiased Score Expansion and application to CMB lensing},
PRD, 105, 103531,
arXiv:2112.09354

\bibitem{CMBLensing}
% CMBLensing.jl — cite the published paper (Table II numbers used for comparison)
Millea, M., et al.\ 2021,
\textit{Optimal CMB Lensing Reconstruction and Parameter Estimation with SPTpol Data},
ApJ, doi:10.3847/1538-4357/ac02bb,
arXiv:2012.01709

\bibitem{Robnik:2022}
Robnik, J., De Luca, G.\ B., Silverstein, E.\ \& Seljak, U.\ 2022,
\textit{Microcanonical Hamiltonian Monte Carlo},
arXiv:2212.08549

\bibitem{Bayer:2023}
Bayer, A.\ E., Seljak, U.\ \& Modi, C.\ 2023,
\textit{Microcanonical Langevin Monte Carlo for field-level inference of CMB},
arXiv:2307.09504

\bibitem{Reinecke:2023}
Reinecke, M., Belkner, S.\ \& Carron, J.\ 2023,
\textit{Faster CMB lensing with ducc0},
arXiv:2304.10431

\bibitem{Millea:2020}
Millea, M., Anderes, E.\ \& Wandelt, B.\ D.\ 2020,
\textit{Bayesian delensing of the CMB},
PRD, 102, 123542,
arXiv:2002.00965

\bibitem{Doeser:2024}
Doeser, L.\ \& Jasche, J.\ 2024,
arXiv:2606.10023

\bibitem{Drouin:2023}
% Score-based / diffusion generative lensing — representative reference
\todo{Fill in the correct diffusion-lensing citation from literature.md}

\end{thebibliography}

\end{document}
